\documentclass[11pt]{article}

\usepackage[a4paper,margin=1in]{geometry}
\usepackage[T1]{fontenc}
\usepackage[utf8]{inputenc}
\usepackage{lmodern}
\usepackage{microtype}
\usepackage{setspace}
\usepackage{parskip}
\usepackage{hyperref}
\usepackage{enumitem}
\usepackage{booktabs}
\usepackage{longtable}
\usepackage{array}

\onehalfspacing

\title{Constitutional Force and Applicability\\
in the Finite-Capacity AI Constitution}
\author{Panagiotis Kalomoirakis\\
Independent Researcher, Synkyria Project}
\date{April 2026}

\begin{document}

\maketitle

\begin{center}
\small \textcopyright\ Panagiotis Kalomoirakis 2026. All rights reserved.
\end{center}

\begin{abstract}
This note clarifies the constitutional force, applicability, and internal stratification of the Finite-Capacity AI Constitution. It does not replace or rewrite the constitutional article draft. Its narrower purpose is to explain how the Constitution should be read canonically: what in it functions as universal minimum constitutional core, what functions as stronger governance-level constitutional discipline, and what functions only where the system class, runtime structure, or domain profile makes richer historical qualification necessary. The note is written in order to prevent two opposite errors: first, the under-reading according to which the Constitution is merely a loose ethical statement; second, the over-reading according to which every AI system would need to instantiate a rich history-sensitive runtime architecture in order to count as constitutional.
\end{abstract}

\section{Status of this note}

This document is an interpretive and classificatory companion to the \emph{Finite-Capacity AI Constitution}. It does not alter the constitutional articles. It clarifies their force, scope, and gradation.

Its purpose is therefore limited:
\begin{quote}
to explain how the Constitution should be read, applied, and eventually stabilised as a canonical object without overburdening the constitutional text itself.
\end{quote}

\section{Why a separate note is needed}

A constitutional article draft must remain concise enough to function as a constitutional object. If it tries to carry, in its own body, every question of applicability, gradation, runtime richness, and canonical transition, it risks losing constitutional clarity.

At the same time, once the constitutional text is taken seriously, legitimate questions arise:
\begin{itemize}[leftmargin=2em]
    \item Which articles are truly universal?
    \item Which articles are stronger governance clauses rather than primitive minima?
    \item Does the Constitution require every AI system to become a rich history-sensitive runtime?
    \item How should the Constitution move from draft status toward canonical status?
\end{itemize}

This note answers those questions without rewriting the constitutional core.

\section{The constitutional problem}

The \emph{Finite-Capacity AI Constitution} states a general structural grammar for accountable AI under bounded conditions. It does not define a single implementation family, a single runtime architecture, or a single audit stack. Its central burden is narrower: it identifies the minimum structural conditions under which constitutional accountability remains possible when capacity is finite.

This immediately implies that the Constitution must be read at more than one level. Some provisions express the bare universal core. Others stabilise that core in stronger public-facing, review-facing, or historically sensitive settings. If these levels are not distinguished, two distortions appear:
\begin{enumerate}[leftmargin=2em,label=\arabic*.]
    \item The Constitution is misread as too weak, because its richer clauses are ignored.
    \item The Constitution is misread as too strong, because every clause is treated as equally primitive and equally universal.
\end{enumerate}

The present note therefore introduces a graded reading.

\section{Three constitutional strata}

The Constitution should be read through three strata of force.

\subsection{Tier A: Universal constitutional minimum}

Tier A contains the articles or clauses that apply to any AI system that wishes to make a serious constitutional claim under finite capacity.

These provisions form the \textbf{bare universal core}. Their purpose is not to create a rich runtime architecture, but to prevent constitutional language from floating free of minimal structure.

Tier A includes, in substance:
\begin{itemize}[leftmargin=2em]
    \item finite capacity,
    \item admissibility before execution,
    \item explicit boundary acts,
    \item lawful holding and refusal,
    \item witness obligation,
    \item reviewability with protected non-disclosure,
    \item declared domain projection and scope,
    \item and review before trust.
\end{itemize}

An AI system that cannot satisfy Tier A should not be treated as constitutionally formed.

\subsection{Tier B: Governance-strengthening constitutional discipline}

Tier B contains provisions that remain constitutional in character, but do not function as equally primitive minima. Their role is to strengthen governance integrity, interpretive discipline, and public accountability once a system makes stronger constitutional claims.

Tier B includes, in substance:
\begin{itemize}[leftmargin=2em]
    \item non-sacrifice and slack preservation,
    \item level integrity and qualification inheritance,
    \item extension acts and lawful re-description,
    \item human authority, responsibility, and redress,
    \item and drift / failure-signature / revision discipline.
\end{itemize}

These are not optional in any trivial sense. Rather, they become increasingly necessary as the system's public claims, review surfaces, deployment sensitivity, or certificate surfaces become stronger.

\subsection{Tier C: Historically sensitive or profile-dependent constitutional qualification}

Tier C contains clauses whose constitutional necessity depends on the class of system, the domain profile, or the runtime architecture. These clauses become necessary where terminal surfaces are not semantically sufficient by themselves.

Tier C includes, in particular:
\begin{itemize}[leftmargin=2em]
    \item historical continuity as qualification,
    \item lineage preservation where terminally similar verdicts may arise from historically non-equivalent trajectories,
    \item and stronger revision lineage where version change would otherwise erase accountable history.
\end{itemize}

Tier C does \emph{not} mean that every AI system must instantiate a rich history-sensitive machine. It means only that where a system produces terminal verdicts whose meaning depends on trajectory, constitutional interpretation must preserve the relevant history rather than collapse everything into one terminal summary.

\section{The Constitution does not require all systems to become TER}

A frequent misreading would be the following:
\begin{quote}
If the Constitution includes qualification inheritance, lineage-sensitive interpretation, and revision lineage, then every AI system must become a rich runtime of the same class as a highly developed history-sensitive architecture.
\end{quote}

This conclusion is incorrect.

The Constitution does not require every AI system to instantiate:
\begin{itemize}[leftmargin=2em]
    \item morphogenetic episodes,
    \item multi-layer closure semantics,
    \item rich lineage triptychs,
    \item or history-sensitive field-level certification.
\end{itemize}

Those belong to richer runtime families and stronger implementation classes.

What the Constitution requires is weaker and more general:
\begin{quote}
where the meaning of an act, verdict, or certificate depends on history rather than terminal surface alone, relevant history must not be erased.
\end{quote}

This principle is constitutional. A particular runtime that gives it a rich operational form is an implementation family, not the Constitution itself.

\section{Why finite capacity comes first}

The order of the constitutional grammar is not arbitrary.

Finite capacity comes first because without it the need for boundary discipline becomes easier to deny. If one assumes unlimited capacity, perfect execution, infinite recoverability, and unconstrained review, then interruption, refusal, slack preservation, and witness discipline all begin to look contingent or secondary.

The Constitution rejects that illusion. It begins from finite capacity because:
\begin{itemize}[leftmargin=2em]
    \item execution is bounded,
    \item review is bounded,
    \item reserve is bounded,
    \item interpretation is bounded,
    \item and lawful continuation cannot be presumed indefinitely.
\end{itemize}

Once finite capacity is acknowledged, admissibility before execution follows naturally. Once admissibility is required, explicit acts become necessary. Once acts matter, witness becomes necessary. Once witness exists, scope and reviewability must be disciplined. The structure is therefore cumulative rather than arbitrary.

\section{Bare core versus full constitutional minimum}

The present note distinguishes between two important formulations.

\subsection{Bare core}

The \textbf{bare core} is the shortest constitutional minimum that can still support a serious claim of bounded accountability:
\begin{itemize}[leftmargin=2em]
    \item finite capacity,
    \item admissibility before execution,
    \item explicit boundary acts,
    \item witness,
    \item reviewability,
    \item declared scope,
    \item and review before trust.
\end{itemize}

\subsection{Full constitutional minimum}

The \textbf{full constitutional minimum} consists of the bare core together with the stabilising governance clauses that protect the core from predictable collapse:
\begin{itemize}[leftmargin=2em]
    \item non-sacrifice,
    \item qualification inheritance,
    \item lawful extension acts,
    \item revision discipline,
    \item and responsibility / redress discipline.
\end{itemize}

The current \emph{Finite-Capacity AI Constitution} should be read as expressing the \textbf{full constitutional minimum}, not merely the bare core.

\section{Applicability rule}

The Constitution applies according to the following rule:

\begin{quote}
Tier A applies to any system making a serious constitutional claim under finite capacity. Tier B applies wherever the system makes stronger governance, review, certificate, or public-assurance claims. Tier C applies wherever terminal surfaces are not semantically sufficient without historical qualification.
\end{quote}

This rule preserves both seriousness and proportionality.

It preserves seriousness because no system may evade the bare minimum by appealing to convenience.

It preserves proportionality because not every system is forced into the richest historical or runtime-sensitive form.

\section{Canonical transition}

The current constitutional text remains explicitly identified as an article draft. That is appropriate at the present stage. A transition from draft status toward canonical status should not occur merely by declaration.

A constitutional object moves toward canonical status when the following conditions are satisfied:

\begin{enumerate}[leftmargin=2em,label=\arabic*.]
    \item the constitutional core is stable in wording and force;
    \item its companion objects exist in sufficiently mature form;
    \item the distinction between universal minimum and richer implementation class is explicit;
    \item the review and evidence burden is no longer carried by implication alone;
    \item and the object can be cited as a constitutional reference point without constant internal apologetic clarification.
\end{enumerate}

In practical terms, this means that the constitutional draft is canonically strengthened when it stands beside:
\begin{itemize}[leftmargin=2em]
    \item an application profile,
    \item an evidence specification,
    \item and a clarification note on limits and misreadings.
\end{itemize}

At that point, the constitutional text may remain unchanged while its canonical authority increases through the surrounding clarified architecture.

\section{Companion instruments}

The Constitution should therefore be read together with a small family of companion instruments:
\begin{itemize}[leftmargin=2em]
    \item \textbf{Application profiles}, which project the constitutional grammar into concrete system classes;
    \item \textbf{Evidence specifications}, which state what review surfaces must exist;
    \item \textbf{Clarification notes}, which prevent interpretive drift and overclaim;
    \item and, where appropriate, \textbf{domain-specific governance profiles}, which attach constitutional structure to institutionally sensitive environments.
\end{itemize}

These instruments do not replace the Constitution. They govern projection, review, and interpretation.

\section{Interpretive caution}

A constitutional object should neither be under-read nor over-read.

It is under-read when it is treated as a vague ethical declaration without structural force.

It is over-read when it is treated as though every article imposed the richest possible runtime discipline on every AI system.

The correct reading lies between those errors:
\begin{quote}
the Constitution states the full constitutional minimum; its application is graded; its richer historical implications arise where system class, domain profile, or runtime architecture makes them necessary.
\end{quote}

\section{Closing statement}

A canonical Constitution for AI under finite capacity should be both serious and proportionate.

It should be serious enough to require a real boundary grammar rather than ethical slogans alone.

It should be proportionate enough not to confuse the universal minimum with the richest runtime instantiation.

This note therefore clarifies the intended canonical reading of the \emph{Finite-Capacity AI Constitution}: a general constitutional object with graded force, universal minimum core, governance-strengthening clauses, and historically sensitive qualifications where terminal surfaces alone are not enough.

\end{document}